\section*{Bab \ref{ch:g9:algebra} --- Aljabar dan Persamaan}

\begin{solution}{exo:g9:algebra:1}
$4(3x - 2) = 12x - 8$.

$(x+6)(x+2) = x^2 + 2x + 6x + 12 = x^2 + 8x + 12$.

$(2x-3)(x+5) = 2x^2 + 10x - 3x - 15 = 2x^2 + 7x - 15$.

$(x-1)(x-9) = x^2 - 9x - x + 9 = x^2 - 10x + 9$.
\end{solution}

\begin{solution}{exo:g9:algebra:2}
$(x+8)^2 = x^2 + 16x + 64$.

$(3x-4)^2 = 9x^2 - 24x + 16$.

$(x+10)(x-10) = x^2 - 100$.

$(5x+2)(5x-2) = 25x^2 - 4$.
\end{solution}

\begin{solution}{exo:g9:algebra:3}
$101^2 = (100+1)^2 = 10000 + 200 + 1 = 10201$.

$59^2 = (60-1)^2 = 3600 - 120 + 1 = 3481$.

$53 \times 47 = (50+3)(50-3) = 2500 - 9 = 2491$.
\end{solution}

\begin{solution}{exo:g9:algebra:4}
$9x + 12 = 3(3x + 4)$.

$x^2 - 64 = (x+8)(x-8)$.

$x^2 + 14x + 49 = (x+7)^2$.

$18x^2 - 6x = 6x(3x - 1)$.
\end{solution}

\begin{solution}{exo:g9:algebra:5}
$4x + 7 = 19$: $4x = 12$, jadi $x = 3$.

$6x - 5 = 2x + 11$: $4x = 16$, jadi $x = 4$.

$5(x-3) = 3x + 1$: $5x - 15 = 3x + 1$, jadi $2x = 16$ dan $x = 8$.
\end{solution}

\begin{solution}{exo:g9:algebra:6}
$(x+4)(3x-12) = 0$: $x = -4$ atau $x = 4$.

$x^2 - 121 = (x-11)(x+11) = 0$: $x = 11$ atau $x = -11$.

$(2x-1)^2 = 0$: satu-satunya penyelesaiannya $x = \frac12$.

$x^2 = 8x$: $x(x - 8) = 0$, jadi $x = 0$ atau $x = 8$.
\end{solution}

\begin{solution}{exo:g9:algebra:7}
\omterm{ex:g1:subtraction:difference}{Selisih} kuadrat dengan $a = 3x + 2$ dan $b = 5$:
\[
E = (3x + 2)^2 - 5^2 = (3x + 2 + 5)(3x + 2 - 5) = (3x + 7)(3x - 3).
\]
Aturan \omterm{def:g2:mult:def}{hasil kali} \omterm{def:g1:counting:zero}{nol}: $3x + 7 = 0$ memberi $x = -\frac73$, dan
$3x - 3 = 0$ memberi $x = 1$.
\end{solution}

\begin{solution}{exo:g9:algebra:8}
$3x - 4 \geq 8$: $3x \geq 12$, jadi $x \geq 4$ (titik penuh di $4$,
arsir ke kanan).

$5 - 2x > 1$: $-2x > -4$, membaginya dengan $-2$ membalikkan tandanya:
$x < 2$ (titik berongga di $2$, arsir ke kiri).

$4(x+1) \leq x - 5$: $4x + 4 \leq x - 5$, jadi $3x \leq -9$, jadi
$x \leq -3$ (titik penuh di $-3$, arsir ke kiri).
\end{solution}

\begin{solution}{exo:g9:algebra:9}
Misalkan $x$ bilangannya: $3x + 7 = 5x - 9$. Lalu $16 = 2x$, jadi
$x = 8$. Periksa: $3 \times 8 + 7 = 31$ dan $5 \times 8 - 9 = 31$.
\end{solution}

\begin{solution}{exo:g9:algebra:10}
\emph{1.} $(n+1)(n-1) = n^2 - 1$ (kesamaan ketiga). Tetangga $n$
adalah $n - 1$ dan $n + 1$, jadi \omterm{def:g2:mult:def}{hasil kalinya} $n^2 - 1$: misalnya
$24 \times 26 = 625 - 1 = 624$.

\emph{2.} Tiga bilangan bulat berurutan dapat ditulis $n - 1$, $n$,
$n + 1$. \omterm{def:g1:addition:def}{Jumlahnya} adalah
\[
(n-1) + n + (n+1) = 3n ,
\]
yaitu \omterm{def:g9:arith:divisor}{kelipatan} $3$ (tiga kali yang di tengahnya).
\end{solution}

\begin{solution}{pb:g9:algebra:1}
\textbf{1.} $(10a + b) - (10b + a) = 9a - 9b = 9(a - b)$: jadi
\omterm{ex:g1:subtraction:difference}{selisih} antara bilangan dua angka dan bilangan baliknya selalu
\omterm{def:g9:arith:divisor}{kelipatan} $9$. Periksa: $72 - 27 = 45 = 9 \times 5$, dan
$a - b = 7 - 2 = 5$.

\textbf{2.} Dengan \omterm{def:g9:algebra:expand}{menjabarkan} ruas kanannya:
$9(11a + b) + a + b + c = 99a + 9b + a + b + c
= 100a + 10b + c$. Jadi bilangannya sama dengan \omterm{def:g9:arith:divisor}{kelipatan} $9$ ditambah
\omterm{def:g1:addition:def}{jumlah} angkanya: kalau dibagi $9$, keduanya meninggalkan \omterm{def:g3:division:remainder}{sisa} yang sama.
Khususnya \omterm{def:g3:division:remainder}{sisa} $0$ untuk yang satu berarti \omterm{def:g3:division:remainder}{sisa} $0$ untuk yang
lain: itulah aturan sembilan, yang terbukti.

\textbf{3.} $9(11a + b)$ juga \omterm{def:g9:arith:divisor}{kelipatan} $3$, jadi kesamaan yang
sama menunjukkan bahwa bilangannya dan \omterm{def:g1:addition:def}{jumlah} angkanya meninggalkan
\omterm{def:g3:division:remainder}{sisa} yang sama pada pembagian dengan $3$: itulah aturan tiga.

\textbf{4.} Tulislah $47 = 9k + 2$ dan $86 = 9l + 5$. Lalu
\[
47 \times 86 = (9k + 2)(9l + 5)
= 9\,(9kl + 5k + 2l) + 10 :
\]
\omterm{def:g2:mult:def}{hasil kalinya} meninggalkan \omterm{def:g3:division:remainder}{sisa} yang sama seperti $2 \times 5 = 10$,
yaitu $1$. Karena itu hasil yang benar pun harus meninggalkan \omterm{def:g3:division:remainder}{sisa}
$1$ --- dan $4\,042$ memang begitu. Tetapi ujiannya hanya melihat
sisanya: $4\,042$ dan $4\,402$ (yang angkanya tertukar) lulus dengan
cara yang sama, jadi lulus tidak membuktikan apa pun --- gagal itulah
yang menghukum.

\textbf{5.} Dengan \omterm{def:g9:algebra:expand}{menjabarkan}: $11(91a + 9b + c) - a + b - c + d
= 1001a + 99b + 11c - a + b - c + d
= 1000a + 100b + 10c + d$. Jadi bilangannya \omterm{def:g6:wholes:divisible}{habis dibagi} $11$
persis ketika \omterm{def:g1:addition:def}{jumlah} bergantinya $-a + b - c + d$ juga demikian.
Untuk $2\,926$: $-2 + 9 - 2 + 6 = 11$, yang \omterm{def:g6:wholes:divisible}{habis dibagi} $11$: jadi
ya, $2\,926 = 11 \times 266$.

\textbf{6.} $742 - 247 = 495$; $495 + 594 = 1\,089$. Contoh yang sah
mana pun mendarat di $1\,089$ juga.

\textbf{7.} $(100a + 10b + c) - (100c + 10b + a) = 99(a - c)$.
Untuk $a - c = 2, 3, \dots, 9$:
\[
198,\ 297,\ 396,\ 495,\ 594,\ 693,\ 792,\ 891 .
\]

\textbf{8.} $100(m - 1) + 90 + (10 - m) = 100m - 100 + 90 + 10
- m = 99m$. Jadi ketiga angka $99m$ adalah $m - 1$, lalu $9$,
lalu $10 - m$ (periksa $495$: $m = 5$ memberi $4$, $9$, $5$).

\textbf{9.} Bilangan balik $99m$ punya angka $10 - m$, $9$,
$m - 1$, jadi nilainya $100(10 - m) + 90 + (m - 1)$. Dengan menjumlahkan:
\[
100\,(m - 1 + 10 - m) + 180 + (10 - m + m - 1)
= 100 \times 9 + 180 + 9 = 1\,089 ,
\]
yang tidak bergantung pada $m$. Berapa pun bilangan awalnya,
pesulapnya aman.

\textbf{10.} Kalau $a = c$, pengurangannya memberi $0$ dan triknya
mati seketika. Kalau $a - c = 1$, pengurangannya memberi $99$, yang
harus diperlakukan sebagai untai tiga angka $099$ (baliknya $990$,
$99 + 990 = 1\,089$: terselamatkan!). Catatan kecilnya: \omterm{ex:g1:subtraction:difference}{selisih}
$1$ berhasil hanya kalau penontonnya menulis hasil pengurangannya
dengan tiga angka, termasuk \omterm{def:g1:counting:zero}{nol} di depannya --- menuntut jarak
sedikitnya $2$ menghindari perdebatannya.

\textbf{11.} $2(5m + 7) + d - 14 = 10m + 14 + d - 14 =
10m + d$: puluhannya (dan ratusannya, untuk Oktober--Desember)
menunjukkan bulannya, satuannya (dan puluhannya) tanggalnya.

\textbf{12.} $2(5m + 9) + d = 10m + 18 + d$: kurangkan $18$.

\textbf{13.} Menurut pertanyaan 2, $n - (\text{jumlah angka } n) =
9(11a + b)$ untuk $n$ tiga angka: jadi selalu \omterm{def:g9:arith:divisor}{kelipatan} $9$
(kesamaan serupa berlaku untuk panjang berapa pun). \omterm{def:g5:division:multiple}{Kelipatan} $9$
punya \omterm{def:g1:addition:def}{jumlah} angka yang juga \omterm{def:g9:arith:divisor}{kelipatan} $9$ (pertanyaan 2 lagi), jadi
angka yang dibacakan penontonnya harus dilengkapi sampai \omterm{def:g9:arith:divisor}{kelipatan}
$9$ berikutnya: angka yang hilang adalah pelengkap itu. Kekaburannya:
kalau angka yang dibacakan sudah berjumlah \omterm{def:g9:arith:divisor}{kelipatan} $9$, angka yang
dicoret dapat $0$ atau $9$ --- mengecualikan $0$ menghapus keraguannya.

\textbf{14.} $(10t + 5)^2 = 100t^2 + 100t + 25 =
100\,t(t + 1) + 25$: tulislah $t(t+1)$, lalu tempelkan $25$ di
belakangnya. $45^2$: $4 \times 5 = 20$, jadi $2\,025$. $85^2$:
$8 \times 9 = 72$, jadi $7\,225$. $105^2$: $10 \times 11 = 110$,
jadi $11\,025$.

\textbf{15.} Lima yang berurutan: $n + (n+1) + (n+2) + (n+3) +
(n+4) = 5n + 10 = 5(n + 2)$: yaitu \omterm{def:g9:arith:divisor}{kelipatan} $5$, selalu. Empat
yang berurutan: $n + (n+1) + (n+2) + (n+3) = 4n + 6 =
4(n + 1) + 2$: yaitu \omterm{def:g3:division:remainder}{sisa} $2$ pada pembagian dengan $4$, jadi tidak
pernah menjadi \omterm{def:g9:arith:divisor}{kelipatannya}. Aljabar mengubah ``tampaknya berhasil'' menjadi
\emph{selalu}, dan ``aku tidak menemukan contohnya'' menjadi
\emph{tidak pernah} --- dua kata yang tidak dapat dicapai pengujian
sebanyak apa pun (\cref{met:g7:literal:test}).

\end{solution}
